\documentclass[12pt]{article}
\usepackage[T1]{fontenc}
\usepackage[utf8]{inputenc}
\usepackage{lmodern}
\usepackage{textcomp}
\usepackage{enumitem}
\usepackage{geometry}
\usepackage{graphicx}
\usepackage{amsmath, amssymb}
\geometry{margin=1in}
\begin{document}
\begin{center}
Sociology - Undergraduate Level Assessment 14 \
Total Marks: 40 \
Answer all questions.
\end{center}

\section*{Multiple Choice (10 marks)}
\begin{enumerate}[label=\arabic*.]
\item Mulholland (1998) argued that privatization changed the relationship between companies and managers in that:
\begin{enumerate}[label=(\alph*)]
    \item flexibility was reduced by the introduction of detailed daily worksheets
    \item the state had greater control than managers over production processes
    \item ownership was transferred from small shareholders to senior managers
    \item employment depended on performance rather than trust and commitment
\end{enumerate}
\item Economic aid has largely failed to promote modernization in the developing countries because:
\begin{enumerate}[label=(\alph*)]
    \item there are no clearly defined projects into which the money can be directed
    \item the United Nations has refused to call on rich countries to provide it
    \item debt repayments with interest can be greater than the amount of money received
    \item debt repayments with interest can be greater than the amount of money received
\end{enumerate}
\item In modern societies, social status is typically measured by a person's:
\begin{enumerate}[label=(\alph*)]
    \item age
    \item income
    \item verbal fluency
    \item occupation
\end{enumerate}
\item Durkheim defined social facts as:
\begin{enumerate}[label=(\alph*)]
    \item ways of acting, thinking, and feeling that are collective and social in origin
    \item the way scientists construct knowledge in a social context
    \item data collected about social phenomena that are proven to be correct
    \item ideas and theories that have no basis in the external, physical world
\end{enumerate}
\item Bourdieu attributed the reproduction of class to:
\begin{enumerate}[label=(\alph*)]
    \item cults of the capital
    \item capital culture
    \item cultural capital
    \item culpable capture
\end{enumerate}
\end{enumerate}

\section*{True / False (10 marks)}
\textit{Answer True or False}
\begin{enumerate}[label=\arabic*.]
\setcounter{enumi}{5}
\item True or False: Dahrendorf, Rex, and Habermas emphasized the significance of power, domination, and conflict in social structures.
\item True or False: Decarceration involves the increased use of imprisonment as the primary method of punishment in society.
\item True or False: The social democratic model is characterized by loyalty to the state and traditional values in welfare provision.
\item True or False: A fixed geographical location is not considered a necessary criterion of community by Fulcher \& Scott.
\item True or False: A document is considered 'authentic' if it is a 'sound' original or a reliable copy of known authorship.
\end{enumerate}

\section*{Long Form Response (20 marks)}
\textit{Answer each question with a clear and complete explanation.}
\begin{enumerate}[label=\arabic*.]
\setcounter{enumi}{10}
\item Discuss the concept of scapegoating in society. How does attributing societal problems to a specific group impact social dynamics and contribute to conflict?
\item Discuss the key processes of decision-making in the USA as identified by Domhoff, and critically analyze why the exploitation process is not included in his framework.
\end{enumerate}

\vfill
\noindent\textit{End of Paper}
\end{document}